\section{Computer-assisted existence results}
\label{sec:certified-existence}

We distinguish the three computer-assisted existence results below from
the numerically observed ladder. The proof establishes a localized mode
at a parameter pair inside each specified box. It does not certify every
observed radiation minimum, an infinite ladder, or a nonlinear
nonradiating solution.

\subsection{Model and precise statement}
\label{sec:proof-model}
Consider the complex Klein--Gordon equation in three spatial dimensions,
\begin{equation}
 \partial_t^2\phi-\Delta\phi+U'(|\phi|^2)\phi=0,
 \qquad U(S)=S-S^2+\tfrac12 S^3.
 \label{eq:proof-kg}
\end{equation}
A stationary Q-ball has the form $\phi_0=e^{i\omega t}f(r)$, where
\begin{equation}
 f''+\frac{2}{r}f'=(1-\omega^2)f-2f^3+\tfrac32 f^5,
 \qquad f(0)=a,\quad f'(0)=0.
 \label{eq:proof-profile}
\end{equation}
We use a real orthonormal basis $Y_{\ell m}$ of spherical harmonics and
write a linear perturbation as
\begin{equation}
 \delta\phi=e^{i\omega t}
 \left[\frac{A(r)}{r}e^{i\rho t}
       +\frac{B(r)}{r}e^{-i\rho t}\right]Y_{\ell m}.
 \label{eq:proof-mode}
\end{equation}
In a complex harmonic basis the second sideband instead contains
$Y_{\ell m}^*$. With $S=f^2$, substitution in the linearization of
\eqref{eq:proof-kg} gives
\begin{equation}
 \begin{aligned}
 A''&=\left[1-4S+\tfrac92 S^2-(\omega+\rho)^2
                  +\frac{\ell(\ell+1)}{r^2}\right]A+CB,\\
 B''&=\left[1-4S+\tfrac92 S^2-(\omega-\rho)^2
                  +\frac{\ell(\ell+1)}{r^2}\right]B+CA,
 \qquad C=-2S+3S^2.
 \end{aligned}
 \label{eq:proof-radial-system}
\end{equation}
Regularity requires $A,B=O(r^{\ell+1})$. Define
\begin{equation}
 \kappa_0=\sqrt{1-\omega^2},\qquad
 k=\sqrt{(\omega+\rho)^2-1},\qquad
 \kappa_c=\sqrt{1-(\omega-\rho)^2}.
 \label{eq:proof-channel-numbers}
\end{equation}

\noindent\textbf{Computer-assisted existence statement.}
For each row of Table~\ref{tab:proof-cases}, the corresponding certificate
encloses a quadruple $(a_*,c_*,\rho_*,\omega_*)$ with $c_*>0$ such that
\eqref{eq:proof-profile} has a positive global solution satisfying
$r e^{\kappa_0r}f(r)\to c_*$, and
\eqref{eq:proof-radial-system} has a nonzero real solution regular at the
origin with
\begin{equation}
 |A(r)|+|B(r)|\leq K_J e^{-\kappa_c r}\quad(r\geq L),
 \qquad e^{\kappa_c r}B(r)\longrightarrow1.
 \label{eq:proof-decay}
\end{equation}
The resulting spatial fields in \eqref{eq:proof-mode} are smooth and
exponentially localized. The certification is
conditional on the correctness of the stated interval-arithmetic
implementation and its underlying arithmetic libraries.

\begin{table}[htbp]
\centering
\small
\begin{tabular}{ccll}
\hline
Case & $\ell$ & Enclosure of $\omega_*^2$ & Enclosure of $\rho_*$ \\
\hline
C0 & 0 & $0.797676787111865191750\pm4.8\,10^{-22}$
       & $1.744617544837398735781\pm6.8\,10^{-22}$ \\
C1 & 0 & $0.6851289044582160933\pm2.5\,10^{-20}$
       & $1.6903565973281434568\pm4.7\,10^{-20}$ \\
C2 & 1 & $0.7544960137826372331\pm7.8\,10^{-20}$
       & $1.826342030181069956\pm3.0\,10^{-19}$ \\
\hline
\end{tabular}
\caption{The three certified parameter boxes, displayed as outward-rounded
decimal enclosures. The proof uses the full exact dyadic boxes in
$(a,c,\rho,\omega)$, not these two-coordinate summaries.
Their numerical ladder
assignments are $(\ell,n)=(0,1),(0,2),(1,1)$; the labels $n$ are not
separately certified nodal counts.}
\label{tab:proof-cases}
\end{table}

For the enclosed parameter pair and its fixed radial angular channel, every
solution of \eqref{eq:proof-radial-system} that is square integrable on
$[L,\infty)$ is a scalar multiple of the constructed solution. This is
\emph{geometric radial simplicity}. It neither establishes algebraic
simplicity of a dynamical generator nor excludes Jordan chains. For C2,
the three-dimensional $\ell=1$ angular multiplet remains degenerate.
The translation mode at $\rho=0$ does not belong to the fixed nonzero
frequency considered here.

The term \emph{embedded} is used operationally: $k^2>0$ and
$\kappa_c^2>0$, so a free asymptotic channel is open, while the constructed
mode is exponentially localized and has no oscillatory radiation tail.
We do not identify the essential spectrum of a fully specified operator
pencil in this proof. The perturbation is periodic after removing
$e^{i\omega t}$; the two laboratory frequencies are $\omega\pm\rho$,
with no claim about their rational commensurability.

\subsection{Infinite-domain construction and inclusion}
\label{sec:proof-tails}
The matching radius $L$ joins an interval-integrated core to an analytic
infinite exterior. The potential and its coupling are retained beyond
$L$. Appendix~\ref{supp:proof-lemmas} gives the complete
analytic lemmas used below, including the origin start and the error
bounds for every Taylor step.

For the profile, $\widetilde u=r e^{\kappa_0r}f$ satisfies
\begin{equation}
 \widetilde u(r)=c+\int_r^\infty
 \frac{1-e^{-2\kappa_0(s-r)}}{2\kappa_0}
 \left(-2f(s)^2+\tfrac32f(s)^4\right)\widetilde u(s)\,\dd s.
 \label{eq:proof-profile-volterra}
\end{equation}
The certified exterior constants give a uniform contraction in a positive
ball about $c$. For the mode, the free open and closed Green kernels
produce a Volterra equation with integrable coefficients. Ordered
iteration has factorial error majorants, and yields a normalized
exponentially decaying Jost solution. The $\ell=1$ construction retains
the exact Riccati functions and the centrifugal term. These constructions
supply continuous corrections to both the profile Cauchy data and the
mode Cauchy data at $L$, uniformly on the parameter box.

\subsection{Finite-dimensional certificate}
\label{sec:proof-krawczyk}
Let $z=(a,c,\rho,\omega)$. The four-component matching map is
$H=H_0+E$. Its first two components match profile value and derivative;
the last two are scaled Wronski forms
$W(\Psi,R_i)=\Psi\cdot R_i'-\Psi'\cdot R_i$, where $R_1,R_2$ are
the two regular origin solutions. The nominal map $H_0$ uses the free
Jost data and $E$ bounds the infinite-tail correction.
For $\ell=0$, these pairings equal the origin values of $\Psi$.
For $\ell=1$, the regular behaviors $r^2e_i$ pair with singular
behaviors $r^{-1}e_i/3$ to give $\delta_{ij}$. Thus their joint vanishing
is precisely regularity of the Jost solution at the origin.

The core integration encloses $H_0(z_0)$ and $DH_0(Z)$ using complex
analytic a priori bounds and Taylor remainders. At the origin, the
profile integral equation has weight mass $r^2/6$; for
$A=r^2\alpha$, $B=r^2\beta$ the $\ell=1$ weight mass is $r^2/10$.
Parameter derivatives are integrated as part of the same enclosed
system. All step choices and preconditioners become exact dyadic input;
acceptance uses directed arithmetic.

Write $Z=z_0+[-\delta,\delta]$, let $D$ enclose $DH_0(Z)$, and let
$|E_i|\leq\epsilon_i$ throughout $Z$. A fixed dyadic matrix $Y$ gives
\begin{equation}
 \begin{split}
 z_0-YH_0(z_0)+(I-YD)[-\delta,\delta]
       +|Y|[-\epsilon,\epsilon]&\subset\operatorname{int}Z,\\
 \max_i\sum_j |I-YD|_{ij}\frac{\delta_j}{\delta_i}&<1.
 \end{split}
 \label{eq:proof-krawczyk-inclusion}
\end{equation}
The first inclusion makes $z\mapsto z-YH(z)$ a continuous self-map of
$Z$. Brouwer gives a fixed point; the weighted matrix inequality makes
$Y$ invertible, so $H(z_*)=0$. This proof requires continuity, but no
derivative bound, for $E$ and does not assert uniqueness of $z_*$.
Positive core enclosures, a minimum-principle argument on the remaining
finite interval, and the positive exterior fixed point establish $f>0$.
Finally $2\kappa_0>\kappa_c$ permits two bounded oscillatory comparison
solutions at infinity. Their nonzero mutual Wronskian and the
four-dimensional symplectic Cauchy-data space imply that every
square-integrable solution is a multiple of the decaying Jost solution.
Appendix~\ref{supp:simplicity} supplies this simplicity proof explicitly.

\subsection{Certificate contents and verification scope}
\label{sec:proof-verification-limits}
The defining certificates use $L=40$ for C0, $L=32$ for C1, and $L=36$
for C2, with 256-bit ball arithmetic. The preserved historical records
identify these three JSON certificates by case, filename and SHA-256
checksum. A distribution copy with redacted execution metadata is a
derivative: the archive manifest lists its checksum separately from the
historical checksum, while retaining the numerical fields. The full
dyadic boxes, channel bounds, tail constants,
inclusion results and integration-step records are the proof inputs.
The fresh primary-certificate reproduction uses the preserved historical
proof-code bytes. The distribution provides separately hashed code
derivatives with redacted comments and one metadata string in the unused
starting-value-search mode; \texttt{CODE-HASH-MAP.tsv} records their
correspondence. These changes leave the certificate execution path unchanged.
A filename recorded by an old driver is resolved through its recorded
hash, rather than assumed to identify a later file with the same name.

Additional runs use larger matching radii and 320-bit arithmetic. C2's
original repeat at $L=40$ failed the inclusion test; a subsequent run
with an enlarged $c$-radius passed after reevaluating all bounds on that
box. This is an adaptive additional certificate. The C2 existence claim
rests on the first successful certificate alone.

For C0, $2\kappa_0>\kappa_c$ was recorded and separately checked but
was absent from the historical aggregate success flag. This documented
flag omission is retained. For preprint 0.51 a separate ball-arithmetic
check reconstructs the exact dyadic parameter box from the reproduced
C0 receipt and verifies $2\kappa_0-\kappa_c>0$ throughout that box,
together with $\rho>\omega$. It binds the checked receipt and the checker
by their checksums; neither the historical driver nor its success flag
is amended.
The C1/C2 drivers include the inequality explicitly. Repeated runs and
code review use the same implementation and do not supply an independent
second implementation. Correctness of Python--FLINT, Arb/Acb and the
recorded execution remains part of the computational trust base.

The arithmetic uses midpoint--radius enclosures
\cite{Johansson2017Arb,PythonFlintDocs,FlintSoftware}; the finite-dimensional
inclusion is of Krawczyk type \cite{Krawczyk1969}, with the explicit
Brouwer argument above for the continuous tail correction. The three
linear existence results leave background stability, algebraic spectral
simplicity and certification of the full empirical ladder open.
